% !TEX root = main.tex







 \section{Broader Impacts} \label{sec-outreach}
\heading{Overview.}
With NIST finalizing post-quantum standards and launching threshold-cryptography initiatives, this project tackles whether these standards will truly hold up in practice. It creates a unique opportunity for the PI to deliver timely input to NIST and industry while cultivating a pipeline of cryptography-literate students from a large public university.

\heading{Standards and industry engagement.}
The PI’s work is uniquely positioned to shape these emerging standards. As a former NIST visiting 
faculty member, the PI has an ongoing relationship with NIST and will contribute  
technical notes to its workshops clarifying, \emph{e.g.}, which hash properties suffice for 
instantiability of candidate schemes. More generally, as the project’s questions align directly with NIST’s 
post-quantum and threshold priorities, its outcomes will provide timely, research-driven guidance. 
Beyond NIST, the PI will engage UMass-affiliated Bitcoin Unlimited as well as collaborators at 
MongoDB and J.P.~Morgan to ensure the project’s industrial impact.





\heading{Training and mentorship.}
The PI has an exceptional record of mentoring undergraduates who advanced to PhD programs in cryptography  
(\emph{e.g.}, Riddhi Ghosal at UCLA, Alice Murphy at Oregon State). Current mentees at UMass include 
one undergraduate and three PhD students. With one of the nation’s largest CS programs, UMass offers 
a large pool of undergraduates who will be engaged through independent studies and summer research tied to the project.

\heading{Integration with teaching.}
A key deliverable will be a graduate seminar on idealized models, with case studies on 
$\mathsf{FO}$ and $\mathsf{Dilithium}$ drawn from this project. These will give students 
first-hand exposure to research questions shaping NIST standards. Lecture materials 
will be made public.

\heading{Outreach.} 
The PI will promote \emph{early exposure to project ideas} through three UMass initiatives.  
(1) In ongoing cooperation with the executive board of the \emph{Cybersecurity Club} (250+ members), he will 
launch a \emph{cryptography subgroup} and lead tutorials on his research and on student-driven topics such as cryptocurrency and cryptography for AI.  
(2) In the \emph{Early Research Scholars Program} (ESRP), which introduces first- and second-year students 
to faculty-led research, he will adapt material from the project into an accessible group project.  
(3) Through the \emph{Office of Community, Outreach, and Organizational Learning} (COOL), which connects to 
regional community-college and high-school students, he will host ``community discussions'' on cryptography’s 
role in protecting personal privacy, with examples from post-quantum migration for secure messaging.

\heading{Synergies with AI.}
Through the \emph{UMass AI Security Lab}, the PI will explore synergies between this project 
and AI security. For example, the PI observed that $\mathsf{OAEP}$-based McEliece encryption 
acts as a \emph{pseudorandom code}~\cite{C:CriGun24}, a property useful for watermarking 
large language models. %He will give a talk at the AISec seminar  on connections between  cryptography and  AI security.

 \section{Prior NSF Support}
 
The PI has not had any NSF support in the last five years.

%More generally, the PI's courses serve as excellent preparation for students who will enter the workforce.
 
 %More generally, I think cryptography should be pushed further down the curriculum chain to not just undergrads but high school students.  This idea is related to the ``CS for All'' initiative.  %In a sense, this is a ``Cryptography for All'' movement where high school kids understand end-to-end encryption and Facebook's ability to see their data.  I am committed to contacting local high schools to arrange lectures and summer internships.

   %(See my attached Letter of Collaboration from Jane Fountain, Department Chair of  Political Science and Public Policy at UMass.)
%\end{itemize}
%\end{itemize}
%Finally, the PI believes an important aspect of his work is collaboration with applied researchers, which he has a track record of~\cite{SIGMOD:MCOKC15,CCS:KKNO16,KKNO17,ZOSZ17}

\iffalse
 to arrange a visit there.  Beyond simply giving talks, the goal of such visits is to foster a productive dialogue about cryptographic standards and  deployment scenarios for the proposed research.   Furthermore, the PI plans integrate his research and educational plan with Georgetown's new Massive Data Institute (MDI), recently established as part of the McCourt School of Public Policy, and the Georgetown University Medical Center (GUMC).  The MDI provides an additional platform to engage those involved in the policy side of privacy concerns in the era of massive data (\emph{e.g.}, regarding what types of cryptographic technologies can and should be used). The GUMC has a large amount of medical data they wish to store ``in the cloud'' in a privacy-preserving way, to which the proposed research has direct applications.  The PI already has a meeting scheduled with  Subha Madhavan,  Director of the Innovation Center for Biomedical Informatics at the GUMC and  Professor of Oncology, in this regard.
\fi
%We also hope to promote cross-disciplinary collaboration with researchers in the areas of algorithms, databases, and networking.

\iffalse
\begin{itemize}
\item With the rise of cloud computing, our society's need for practical technologies that allow flexible data sharing/processing while offering some degree of data privacy becomes more acute (especially for medical data).  The component of the proposal on encryption that leaks some information addresses the balance of functionality and privacy needed in such applications.   
\item  e protocol designers  with a toolkit of cryptographic primitives more suitable for deployment.
\item 
 remains a difficult problem in practice.  The component of the proposal on program obfuscation has the potential to such technologies that are not only practical but offer provable security guarantees against pirates.
\end{itemize}
\fi

%  A key focus of our integrated outreach plan will be on \emph{industry and policy outreach}, for which Georgetown's location in Washington DC offers an ideal opportunity. 

%As our national security also  depends on cultivating interest in STEM research  as early in education  as possible, we will also \textbf{involve undergraduate and high school students} as appropriate. 
%The PI is currently working with one Georgetown undergraduate (Steven Burchfield) on program obfuscation in practice (Component III-B) and was recently awarded an REU supplement grant on this topic. We are committed to targeting women and minority students for such positions.
%Georgetown is  also unique  in that many of our CS undergraduates  double major in a social science  and go into policy.  We  hope to help such students  gain a deeper understanding of the research side of science, so that they can more effectively influence  policy.  
%This coming year, the PI also plans to help teach a module at a local high school about privacy and the basics of   cryptography,  motivated by sensitive text messages and Facebook posts.  

\iffalse
As our national security also depends on cultivating interest in STEM as early as possible in the educational process, the PI will also reach out to  high school and undergraduate students.     In particular, the PI will  offer \textbf{internships to students at Thomas Jefferson High School} (a nationally-renowned  STEM school that runs an internship program with local universities), as well as participate during the year  in outreach events (such as teaching short modules introducing basic cryptography) at  DC public schools that the Georgetown University Center for Social Justice already has relationships with.  Indeed, the PI believes some such students may have  their interest sparked by more security-oriented topics, and exposure to such material early on is unfortunately rare.  He has also recently secured  REU funding for an undergraduate student through the Security and Software Engineering Research Center at Georgetown, an NSF industry/university cooperative research center.
  He is committed to targeting women and other under-represented  minorities for such positions.
Moreover, Georgetown is unique in that many CS undergraduates double major in a social science and go into policy in some respect.  He hopes to work with such undergraduates so that they may better understand scientific research and therefore more effectively influence policy.
]\fi

%Finally, this proposal will help  establish UMass and western Massachusetts in general in the cryptography community.  %ßThe PI is currently working to significantly  revamp the cryptography curriculum at Georgetown  and build a strong research group (see~\secref{sec-edu}).

 % In order to  solidify the  community and attract  outside attention, the PI   plans to help co-organize a regular ``Western Massachusetts Crypto Day.''  We will also try to establish strategic alliances with the Boston and New York City cryptography research communities.

%Finally, it is worth noting that both   % The PI also plans to hold a weekly or bi-weekly reading group and will continue inviting  high-profile speakers for the departmental colloquium.
   %In addition to disseminating results, our hope is to establish Georgetown and the greater DC area as recognized leaders in the field.

\iffalse 
\section{Prior Work}

Prior work has looked at some cases of trading ``relaxed'' security guarantees for better functionality or efficiency. However, this work has been limited in terms of both applications and the specific relaxations considered, whereas our proposal will fully develop this idea across multiple types of cryptographic schemes and provide a comprehensive treatment.
Below we survey some prior work at a high level. 


\paragraph*{\large{Work related to the component on encryption.}}
Early work in the database community~\cite{} proposed various heuristic encryption schemes  that leak  information in order to allow efficient query processing. 
On the other hand, early work in the cryptographic community met strong security models but required linear-time search~\cite{}.
To the best of our knowledge, the first works to bridge the gap and  treat the provable security of encryption schemes  that leak information were done the PI and his collaborators in the cases of deterministic~\cite{} and order-preserving encryption~\cite{}.   
However, the above schemes are only for \emph{relational} databases.
Several recent works in the database community address the case of nearest-neighbor search  (which the PI plans to address in detail)  but are limited or heuristic in nature~\cite{}.
  %The proposed research will substantially generalize and broaden the applicability of these results.   In particular,  a topic that the PI plans to address in detail.  
%The approach of designing encryption schemes that leak partial information should be contrasted with homomorphic encryption~\cite{}, which  allows processing on encrypted data with strong security guarantees but necessitates specialized and inefficient processing methods.  
We also mention that a large body of recent work in the cryptographic community  considers ``symmetric searchable encryption'' protocols that leak  partial information~\cite{} such as database size and access pattern, but they encrypt the database ``all at once'' rather than on the fly and necessitate specialized implementations.   %In the context of general two-party computation,  %Franklin and Mohassel introduced an influential security model allowing one bit of leakage~\cite{} and showed it leads to more efficient protocols.
%  As we mentioned, our goal is not to treat the  ``relaxed'' security of such higher-level protocols in detail but rather that of the lower-level building blocks like encryption and signature schemes, which can be used to construct them.

\paragraph*{\large{Work related to the component on authentication.}}


We also mention that  a line of work started by Dwork and Naor~\cite{} examines ``proofs of work'' functions that are moderately hard to compute and are inherently sequential (\emph{i.e.}, computing such functions $n$ times takes $n$ times more work), but may be faster to compute in bulk via some trapdoor.  Our concept of moderately-unforgeable signatures is superficially similar but the signing key in our case allows much faster computation of even a \emph{single} signature, and we seek realizations where the signatures are much faster to compute than any known signatures that are existentially unforgeable.  

\paragraph*{\large{Work related to the component on obfuscation.}}
As mentioned above, program obfuscation methods satisfying strong, general-purpose security models have been constructed in recent breakthrough but are very inefficient~\cite{}.   Our inspiration for considering relaxed security models comes from applied work using heuristic but much more practical techniques.  In particular, the goal of the obfuscating a random key within a program is inspired by common techniques used to obfuscate polymorphic viruses and malware~\cite{}, which  embeds a key in an obfuscated program along with encrypted code --- the latter is then decrypted and executed as the ``actual code'' at run time.
The goal of obfuscating conditional branches is inspired by  work by Collberg on ``opaque predicates''~\cite{}.
\fi

\iffalse
\subsection{Evaluation}

Success of the proposed research can be evaluated by number of publications it yields;  we anticipate an average of 1.5  per year. 
It should be noted we are not only hoping to publishing theory-oriented papers, but also more applied papers actually integrating our proposed schemes with applications.
The broader impact can evaluated by number of outreach visits or guest appointments to government/standards bodies; we anticipate about 3 outreach visits or guest appointments per year. 
The educational component can be evaluated by teaching evaluations (our  goal is at least a 4/5 rating for each course; our ratings last year were 4.04 and 4.17) and number of students supervised; we hope to advise an average of two  graduate students, one post-doc, and one undergraduate  per year.  
\fi

\section{Relevance to Secure and Trustworthy Cyberspace}

This project advances SaTC’s core goal of ``sound mathematical and scientific foundations'' by resolving tightness and instantiability for standardized schemes --- an urgent need for post-quantum standards and Internet security.
%The keywords selected for the proposal are ``cryptography; theory'' as the work, while having practical implications, has to do with our theoretical understanding. 

%\section{Prior NSF Support} \label{sec-accom}

%\paragraph*{Prior Accomplishments.}
 % The PI has a track record in the study of idealized models~\cite{C:KilOneSmi10,TCC:GoyOneRao11,EC:LewOneSmi13,PKC:CaoOneZah20,AC:MurONeZah22}.   He also has strong interests in advanced forms of encryption and signatures schemes, such as searchable~\cite{C:BelBolOne07,DBSec:AmaBolOne07,C:BolFehOne08,C:BFOR08,TCC:FulOneRey12,EC:BCLO09,C:BolCheOne11,CCS:KKNO16,HO18,KKNO17,ZOSZ17,AC:CLOZZ18,} and functional encryption~\cite{EPRINT:Oneill10,C:DIJOPP13, CANS:BelONe13,AC:GoyJaiOne16,ITCS:ACFGOT20}.  Finally, he has as also worked  on many other topics in cryptography including    aggregate signatures~\cite{CCS:BGOY07,AC:GLOW12,PKC:GenOneRey18}, deniable encryption~\cite{C:OnePeiWat11,PKC:DecIovOne16},  chosen-ciphertext security~\cite{EC:KilMohOne10,PKC:DFMO14}, leakage resilience~\cite{PKC:DGLOZ16,CANS:BelOneSte17,EPRINT:CKOP20}, and selective-opening security~\cite{AC:HKOZ16,EPRINT:ZahOne20}. 

%\paragraph*{NSF Support and Results.}  
 % Previous support was:
%``Guaranteed-Secure and Searchable Genomic Data Repositories'', Principal Investigator (Co-PI: Ophir Frieder), SaTC NSF EAGER \#CNS-1650419, 2016 - 2019, \$99,999.
%\textbf{Intellectual merit:} Introduced a new concept of ``ad hoc'' multi-input functional encryption~\cite{ITCS:ACFGOT20}. \textbf{Broader impacts:} Gave presentation at DC area ACM Women's Chapter annual meeting, one female REU student.


%\paragraph*{Synergistic Activities.} 
%The PI served on  program committees for Inscrypt 2009, CANS 2010, PKC 2012, Pairing 2012, Eurocrypt 2014, CANS 2014, SCN 2014, SIGIR 2014 Workshop on
%Privacy-Preserving Information Retrieß??val, PKC 2015.  The PI  served on an NSF/SATC small crypto panel (May 2014).  
%The PI spoke at the Boston Freedom in Online Communication Day in March 2013 and speaks regularly at the Security and Software Engineering Research Center (S$^2$ERC) at Georgetown.  For June 2013 %the PI visited Tel-Aviv University and for July 2014 will visit IBM T.J.~Watson.


%The topic of the REU supplement grant overlaps somewhat with~\secref{sec-obf} of the current proposal, as the goal is to attract an undergraduate student to work on this component of the research.